\section{Future directions}\label{sec:future}
% ===============================================================

\textbf{Higher genus and Jacobians.} On a compact Riemann surface $X$ of genus
$g \geq 1$ the equality $\Prin(X) = \ker(\deg)$ fails, and the degree-zero class
group is the Jacobian $\operatorname{Jac}(X) \cong \CC^g/\Lambda$, a complex
torus. Abel's theorem characterises which degree-zero divisors are principal in
terms of the vanishing of a period integral, and Jacobi's inversion theorem
identifies the image; together they are the correct generalisation of
Theorem~\ref{thm:main}. See Miranda \cite[Ch.~VIII]{Miranda}.

\textbf{Riemann--Roch.} Theorem~\ref{thm:main} may be read as a statement about
the dimension $\ell(D)$ of the space of global sections of $\Ostr(D)$: on $\PP$
one has $\ell(D) = \deg D + 1$ for $\deg D \geq 0$ and $\ell(D) = 0$ otherwise.
The Riemann--Roch theorem, $\ell(D) - \ell(K-D) = \deg D + 1 - g$, contains this
as the case $g = 0$ and is the natural next result to study.

\textbf{Vector bundles and cohomology.} Grothendieck's splitting theorem states
that every holomorphic vector bundle on $\PP$ is a direct sum
$\Ostr(n_1)\oplus\cdots\oplus\Ostr(n_r)$; Theorem~\ref{thm:pic} is the case
$r = 1$. In the sheaf-theoretic language, $\Pic(X) \cong H^1(X, \Ostr^\ast)$, and
Lemma~\ref{lem:cousin} is the computation of that cohomology group for $\PP$ by
hand using the two-set cover.

\phantomsection
\addcontentsline{toc}{section}{References}
